\section{第12章 ソフトウェア品質}
{\Huge\bfseries\sffamily 12 章 ソフトウェア品質\par}
{\Large Software Quality の日本語まとめ\par}
{\normalsize SWEBOK Guide v4.0a Chapter 12 を初学者向けに再構成\par}
{\normalsize 作成日：2026 年 5 月 12 日\par}

\begin{pointbox}{この資料の主張}
この資料の主張\par
\term{ソフトウェア品質}{Software Quality}は、単に「バグが少ないこと」ではない。品質とは、利害関係者のニーズを正しく表した要求に、ソフトウェア製品・作業成果物・開発プロセスがどの程度適合しているかを扱う概念である。したがって、品質はテスト工程だけで作り込むものではなく、要求、設計、実装、レビュー、構成管理、測定、是正処置、運用までを通じて管理する対象である。
\end{pointbox}

\begin{pointbox}{この資料の主張}
この資料の読み方\par
専門用語は原則として日本語で示し、必要に応じて英語を括弧内に併記する。英語名は、元資料やツール上の名称を確認したいときの手がかりとして残す。初学者が前提知識なしに読めるように、各節では「何のための概念か」「実務では何を確認すればよいか」を重視する。文体は、だである調に統一する。
\end{pointbox}
